\documentclass[11pt, a4paper, twoside]{article}

% --- Packages ---
\usepackage[utf8]{inputenc}
\usepackage{amsmath, amssymb, amsthm}
\usepackage{geometry}
\usepackage{tcolorbox}
\usepackage{tikz-cd}
\usepackage{fancyhdr}

% --- Page Layout & Styling ---
\geometry{
	top=1in, 
	bottom=1in, 
	left=1.25in, 
	right=1.25in
}

\pagestyle{fancy}
\fancyhf{}
\fancyhead[LE,RO]{\thepage}
\fancyhead[LO]{\small \textit{Lecture Notes: Linear Block Codes}}
\fancyhead[RE]{\small \textit{A Coding Theoretical Interpretation of the SES}}
\renewcommand{\headrulewidth}{0.4pt}

% Custom boxes
\newtcolorbox{conceptbox}[1][]{
	colback=blue!5!white, colframe=blue!75!black,
	fonttitle=\bfseries, title=#1, arc=3pt, boxrule=1pt,
	left=10pt, right=10pt, top=10pt, bottom=10pt
}
\newtcolorbox{mathbox}[1][]{
	colback=gray!5!white, colframe=gray!75!black,
	fonttitle=\bfseries, title=#1, arc=3pt, boxrule=1pt,
	left=10pt, right=10pt, top=10pt, bottom=10pt
}

% --- Macros ---
\newcommand{\F}{\mathbb{F}}
\newcommand{\C}{\mathcal{C}}
\newcommand{\im}{\text{im}}
\newcommand{\kerop}{\text{ker}}

\begin{document}
	
	\begin{center}
		{\Large \textbf{The Physical Meaning of Exactness}} \\[0.5em]
		{\normalsize \textit{A Coding Theoretical Interpretation of the Short Exact Sequence}}
	\end{center}
	\vspace{1em}
	
	While the Short Exact Sequence (SES) is an object of abstract homological algebra, in coding theory, it physically models the transmission of data across a noisy channel. The sequence is defined as:
	
	\[
	\begin{tikzcd}[column sep=large]
		0 \arrow[r] & \F_q^k \arrow[r, "G"] & \F_q^n \arrow[r, "H"] & \F_q^{n-k} \arrow[r] & 0
	\end{tikzcd}
	\]
	
	We analyze the coding-theoretical significance of exactness at each node in this sequence.
	
	\vspace{1em}
	
	\begin{conceptbox}[1. Left Exactness: Unambiguous Information]
		\textbf{Mathematical Condition:} Exactness at $\F_q^k$ requires that $\kerop(G) = \{0\}$. Therefore, $G$ is strictly injective.
		
		\textbf{Coding Interpretation:} 
		The space $\F_q^k$ represents the raw information source. If $G$ were not injective, there would exist distinct messages $m_1 \neq m_2$ such that $m_1 G = m_2 G$. If the receiver obtains this codeword, it would be fundamentally impossible to determine which message was sent. Exactness guarantees that redundancy is added without destroying the uniqueness of the information.
	\end{conceptbox}
	
	\vspace{1em}
	
	\begin{conceptbox}[2. Middle Exactness: The Channel and Error Detection]
		\textbf{Mathematical Condition:} Exactness at $\F_q^n$ requires that $\im(G) = \kerop(H)$.
		
		\textbf{Coding Interpretation:} 
		The ambient space $\F_q^n$ represents the output of the noisy channel. A received word is defined as $y = c + e$, where $c \in \im(G)$ is the transmitted codeword and $e \in \F_q^n$ is the error vector.
		
		Because the sequence is exact, $H$ annihilates any valid codeword ($Hc^T = \mathbf{0}$). When the receiver applies the parity-check matrix to the received word, by linearity:
		$$ s = Hy^T = H(c^T + e^T) = Hc^T + He^T = \mathbf{0} + He^T = He^T $$
		Exactness strips away the message entirely, leaving a \textbf{syndrome} $s$ that depends exclusively on the channel noise. The condition $s = \mathbf{0}$ acts as the mathematical threshold for error detection.
	\end{conceptbox}
	
	\vspace{1em}
	
	\begin{conceptbox}[3. Right Exactness: Cosets and Syndrome Decoding]
		\textbf{Mathematical Condition:} Exactness at $\F_q^{n-k}$ requires that $\im(H) = \F_q^{n-k}$. Therefore, $H$ is surjective.
		
		\textbf{Coding Interpretation:} 
		By the First Isomorphism Theorem, the surjectivity of $H$ implies:
		$$ \F_q^n / \im(G) \cong \F_q^{n-k} $$
		In coding theory, the quotient space $\F_q^n / \C$ defines the \textbf{Standard Array}. The space of all received words $\F_q^n$ is partitioned into exactly $q^{n-k}$ disjoint cosets. 
		
		Each coset takes the form $e + \C$, representing a specific class of error patterns. Right exactness ensures there is a perfect one-to-one correspondence between the $q^{n-k}$ possible hardware syndromes and these error cosets. This is the theoretical foundation of Syndrome Decoding: the receiver computes $s$, uses the isomorphism to identify the corresponding coset, and subtracts the minimum-weight error (the coset leader) to recover $c$.
	\end{conceptbox}
	
\end{document}